\sheet{0}{Getting Started: Wokwi and JSLinux}{Preface, Chapter~1.1}

\begin{goals}
\begin{itemize}
  \item Run, modify and share an Arduino project in the Wokwi simulator.
  \item \emph{Bonus (optional):} compile and run C programs in JSLinux, plot
    results in the terminal, and transfer measured data from Wokwi (Serial
    Monitor) to JSLinux.
  \item Get to know the course library \codefile{common/engine\_signals.h}, the
    \emph{virtual engine} used on all following sheets.
\end{itemize}
\end{goals}

\section*{The virtual engine}

Throughout the exercises we cannot measure on a real engine. Instead, a
\emph{virtual engine} generates the sensor signals introduced in Chapter~1 of
the script. Exactly the same C model is used in three places:

\begin{center}
\small
\begin{tabularx}{\textwidth}{@{}lX@{}}
\toprule
Where & How \\
\midrule
JSLinux & \verb|#include "engine_signals.h"| in your C program \\
Wokwi (sketch) & add \texttt{engine\_signals.h} as an extra file to the project \\
Wokwi (custom chip) & part \texttt{chip-engine-sim}: outputs CRANK, CAM, KNOCK, ION, and
  the ground-truth signals TDC and KNK \\
\bottomrule
\end{tabularx}
\end{center}

The conventions of the model are fixed for all sheets:
\begin{itemize}
  \item 4-cylinder 4-stroke engine, crank angle $\theta\in[0^\circ,720^\circ)$,
    firing order 1--3--4--2, combustion TDC at $120^\circ$, $300^\circ$,
    $480^\circ$, $660^\circ$.
  \item 60-2 trigger wheel: tooth $k=0,\dots,57$ is HIGH for
    $\theta \bmod 360^\circ\in[6k^\circ,6k^\circ+3^\circ)$; teeth 58 and 59 are missing.
  \item Camshaft signal HIGH for $0^\circ\le\theta<180^\circ$.
  \item Knock: damped oscillation with modes at \SI{6.5}{kHz} and
    \SI{10.5}{kHz}, decay time \SI{0.8}{ms}, onset 5--25$^\circ$ after TDC.
  \item Disturbances: broadband noise, valve impacts (\SI{8}{kHz},
    $90^\circ$ after TDC), low-frequency combustion vibration, a
    1\,\% speed ripple at twice crankshaft frequency.
\end{itemize}

%-------------------------------------------------------------------------
\begin{exercise}{First steps in Wokwi}{\Wokwi}
Open \url{https://wokwi.com}, create a new \emph{Arduino Uno} project and replace
\texttt{sketch.ino} and \texttt{diagram.json} by the files in
\codefile{wokwi/sheet00\_hello/}.
\begin{tasks}
  \item Start the simulation, turn the potentiometer and watch the Serial
    Monitor. Which ADC value corresponds to \SI{2.5}{V}? What is the
    voltage resolution (1 LSB) of the 10-bit ADC?
  \item The sketch toggles pin 13 at every sample. Stop the simulation after
    about one second, download the logic analyzer recording (\texttt{.vcd})
    and determine the sampling period. Does it match \texttt{TS\_US}?
  \item Change the sampling rate to \SI{1}{kHz}. Explain why the printed
    \texttt{t\_us} values still show regular increments although the
    \texttt{Serial.print} calls take time.
  \item Wokwi runs in \emph{simulated time}. Compare the simulated time in the
    status bar with the time on your wall clock. Why is this important when
    we later measure execution times?
\end{tasks}
\end{exercise}

\begin{solution}
a) $V = \text{adc}\cdot 5/1023$, so \SI{2.5}{V} $\approx$ 511/512.
1 LSB $= 5/1023\,\text{V} \approx \SI{4.9}{mV}$.

b) Pin 13 toggles at every sample, i.e.\ the period of the rectangular signal on D0
is $2T_s = \SI{20}{ms}$; consecutive edges are $T_s=\SI{10}{ms}$ apart.

c) The loop waits for an absolute deadline \texttt{t\_next} that is incremented by
\texttt{TS\_US}. The time spent in printing only shortens the waiting time, it does
not shift the next deadline (no drift), as long as one loop iteration takes less
than $T_s$. At \SI{1}{kHz} and 115200 baud a line of about 20 characters takes
$\approx\SI{1.7}{ms}$ -- longer than $T_s$! Then the deadlines are missed and
the samples come out late; the \texttt{t\_us} increments become irregular. Remedy:
fewer characters, higher baud rate, or buffer samples and print blocks.

d) The simulator may run faster or slower than real time. \texttt{micros()} and
the logic analyzer use simulated time, so timing \emph{inside} the simulation is
consistent. Wokwi is not cycle-accurate for all peripherals, therefore measured
execution times are only indicative; we count operations instead.
\end{solution}

%-------------------------------------------------------------------------
\begin{exercise}{The engine simulator chip}{\Wokwi}
Every Wokwi project of the following sheets contains the custom chip
\texttt{engine-sim}. Add the files \texttt{engine-sim.chip.c} and
\texttt{engine-sim.chip.json} from \codefile{wokwi/chips/} to the Uno project of
Exercise~0.1 and add the part
\verb|{ "type": "chip-engine-sim", "id": "eng", "attrs": { "vmax": "5", "offset": "2.5" } }|
to \texttt{diagram.json}.
\begin{tasks}
  \item Connect \texttt{eng:CRANK} to logic analyzer D1, \texttt{eng:CAM} to D2
    and \texttt{eng:TDC} to D3 (and \texttt{eng:GND} to ground). Record
    \SI{100}{ms} at \SI{3000}{rpm} and view the \texttt{.vcd} file (e.g.\ in
    PulseView or GTKWave). Identify the missing-tooth gap and verify the
    tooth period.
  \item How many crank teeth lie between the gap and the first TDC pulse?
    Compare with the conventions above.
  \item Change the engine speed with the slider of the chip and repeat the
    tooth-period measurement at \SI{6000}{rpm}.
\end{tasks}
\end{exercise}

\begin{solution}
a) At \SI{3000}{rpm} one revolution takes \SI{20}{ms}; the tooth period
(rising to rising edge) is $\SI{20}{ms}/60=\SI{333}{\micro s}$, across the gap
$3\cdot\SI{333}{\micro s}=\SI{1}{ms}$.

b) The reference edge (first tooth after the gap) is at $\theta=0^\circ$, TDC of
cylinder 1 at $120^\circ$, i.e.\ 20 teeth after the reference edge.

c) At \SI{6000}{rpm}: \SI{167}{\micro s} per tooth, \SI{500}{\micro s} across
the gap.
\end{solution}

\begin{deliverables}
Screenshot of the Wokwi logic analyzer view (or VCD plot) with the gap marked.
Bonus (optional): terminal output of \texttt{hello\_signals}; the CSV-plot
program of 0.B1\,d).
\end{deliverables}

%=========================================================================
\section*{Bonus: JSLinux (optional)}
\begingroup
\setcounter{exercise}{0}\renewcommand{\theexercise}{B\arabic{exercise}}
The lab exercises above use only Wokwi. Interested students can in addition
process the engine signals in plain C on a Linux system that runs in the
browser. The JSLinux exercises of the following sheets build on this one.

\begin{exercise}{First steps in JSLinux}{\JSLinux}
Open \url{https://bellard.org/jslinux} and start the \emph{x86 Alpine Linux}
(console) virtual machine.
\begin{tasks}
  \item Check the toolchain: \texttt{gcc --version} (or \texttt{tcc -v}). Upload
    \texttt{engine\_signals.h}, \texttt{asciiplot.h}, \texttt{csvio.h}
    (\codefile{common/}) and \codefile{bonus/jslinux/sheet00/hello\_signals.c} with
    the upload button below the terminal.
\begin{lstlisting}[style=shell]
gcc -O2 -o hello_signals hello_signals.c -lm
./hello_signals
\end{lstlisting}
  \item Explain the plot of the knock signal: where does the burst start, how
    long does it last, and what is visible before the burst?
  \item Modify the program such that it plots one full engine cycle
    ($720^\circ$) at \SI{3000}{rpm} with \texttt{knock\_prob = 0}. How many
    samples are needed at $f_s=\SI{50}{kHz}$? Identify the valve impacts and
    the combustion vibration in the plot.
  \item Copy 2--3 seconds of the Wokwi Serial Monitor output from
    Exercise~0.1 into a file \texttt{pot.csv} in JSLinux
    (\texttt{cat > pot.csv}, paste, \texttt{Ctrl-D}). Write a short program
    that reads it with \texttt{csv\_read()} and plots the voltage column.
\end{tasks}
\end{exercise}

\begin{solution}
b) The program starts the record at $\theta=125^\circ$, i.e.\ $5^\circ$ after
TDC of cylinder 1. The knock onset lies between $5^\circ$ and $25^\circ$ ATDC;
at \SI{3000}{rpm} ($\SI{18000}{\degree/s}$) $1^\circ$ corresponds to
\SI{55.6}{\micro s}, so the burst appears within the first \SI{1.1}{ms}.
It decays with $\tau=\SI{0.8}{ms}$ and is practically gone after
$3\tau\approx\SI{2.4}{ms}$. Before the burst only noise and the slow
combustion vibration (a half-sine over $60^\circ$) are visible.

c) One cycle lasts $2\cdot 60/3000\,\text{s} = \SI{40}{ms}$, i.e.\ 2000 samples.
Valve impacts appear $90^\circ$ after each TDC (\SI{8}{kHz} bursts at
$\theta = 210^\circ, 390^\circ, 570^\circ, 30^\circ$), the combustion vibration as
broad positive humps directly after each TDC.

d) Example:
\begin{lstlisting}
csv_t d;
if (csv_read("pot.csv", &d) == 0) {
    double *v = csv_col(&d, 3);         /* column "volt" */
    ap_plot(v, d.rows, "potentiometer voltage");
    free(v); csv_free(&d);
}
\end{lstlisting}
\end{solution}
\endgroup
